% !TeX root = ../main.tex


\section{Proposed Method: VLM2Rec}
\label{sec:methods}

To address the SFT-amplified modality collapse (Sec.~\ref{sec:empirical}), we propose \textbf{VLM2Rec} (Fig.~\ref{fig:method}), which trains a VLM with two objectives, without changing its architecture, to recommend the next item by sequence--item similarity.
We follow the modality embedding setting of Sec.~\ref{sec:prompt_templates}.
$\mathcal{L}_{\text{WPCL}}$ strengthens each user's weaker modality (Sec.~\ref{sec:wpcl}), and $\mathcal{L}_{\text{CRTR}}$ softly aligns each user's relative similarity structure over the same items in the text and vision spaces (Sec.~\ref{sec:crtr}); the final objective combines them (Sec.~\ref{sec:final_objective}).
We also propose $\text{VLM2Rec}_{\text{plug}}$ for efficient serving; it transfers the trained item embeddings to SR backbones such as SASRec~\cite{sasrec}, so its latency matches that of conventional SR models (Sec.~\ref{sec:deployment}).

\subsection{Weak-modality Penalized Contrastive Learning}
\label{sec:wpcl}
We propose \textit{Weak-modality Penalized Contrastive Learning} ($\mathcal{L}_{\text{WPCL}}$), which multiplies each user's negative term in $\mathcal{L}_{\text{SFT}}$ (Eq.~\ref{eq:sft_loss}) by a penalty weight that grows as the user's weaker modality lags behind the stronger one.
% It targets the imbalance in Sec.~\ref{sec:empirical}: SFT reduces the loss mainly through the strong modality as a shortcut (A2) and leaves the weak modality barely separating next items from negatives (A3).
$\mathcal{L}_{\text{WPCL}}$ targets the imbalance in Sec.~\ref{sec:empirical}: SFT reduces the loss mainly through the strong modality as a shortcut (A2), so the weak modality barely separates next items from negatives (A3).

\subsubsection{\textbf{User-adaptive Modality Gap Estimation}}
We first measure how clearly each modality separates a user's next item from negatives, as a per-user counterpart of the separability $S^m$ (Sec.~\ref{sec:empirical}).
For user $u$ and modality $m\in\{T, V\}$, the \textit{Discriminative Margin} $\mathcal{M}_{u,m}$ is the similarity of the sequence embedding $\mathbf{e}_u^m$ to the next item minus the mean similarity to the negatives $i^-\in\mathcal{N}_u$:
\begin{equation}
    \mathcal{M}_{u,m} = \text{s}(\mathbf{e}^m_u, \mathbf{e}^m_{i^+}) - \frac{1}{|\mathcal{N}_u|} \sum_{i^- \in \mathcal{N}_u} \text{s}(\mathbf{e}^m_u, \mathbf{e}^m_{i^-}),
\end{equation}
where $\mathcal{N}_u$ holds the next items of the other users in the batch (Sec.~\ref{sec:exp_setup}).
% A larger $\mathcal{M}_{u,m}$, a per-user counterpart of the separability $S^m$ (Sec.~\ref{sec:empirical}), means that modality $m$ separates $u$'s next item more clearly.
At each training step, user $u$'s strong and weak modalities are those with the larger and smaller margins:
\begin{equation}
    \mathcal{M}_{u,\text{strong}} = \max_{m}(\mathcal{M}_{u,m}), \quad \mathcal{M}_{u,\text{weak}} = \min_{m}(\mathcal{M}_{u,m}).
\end{equation}
We identify the weak modality per user because it varies across users (App.~\ref{app:anchor_weak_split}).
$\mathcal{L}_{\text{WPCL}}$ can thus strengthen text for some users and vision for others, unlike gradient modulation methods~\cite{modality_balancing_OGM-GE_1, modality_balancing_AGM_2} based on statistics aggregated over samples (Sec.~\ref{sec:related}).

We define the \textit{Modality Gap} $\Delta_{u,\mathrm{gap}}$ as the difference between these margins:
\begin{equation}
    \Delta_{u,\mathrm{gap}} = \mathrm{sg}[\mathcal{M}_{u,\text{strong}}] - \mathcal{M}_{u,\text{weak}},
    \label{eq:modality_gap}
\end{equation}
where $\mathrm{sg}[\cdot]$ is the stop-gradient operator.
It treats $\mathcal{M}_{u,\text{strong}}$ as a fixed target inside $\Delta_{u,\mathrm{gap}}$ only, so the penalty below can shrink the gap by raising $\mathcal{M}_{u,\text{weak}}$, not by lowering $\mathcal{M}_{u,\text{strong}}$.



\subsubsection{\textbf{Dynamic Penalty from Modality Gap}}
We convert $\Delta_{u,\mathrm{gap}}$ into a penalty weight $w_{u,\text{pen}}$:
\begin{equation}
\label{eq:wpen}
w_{u,\text{pen}}
=
1+\mathrm{Softplus}(\Delta_{u,\mathrm{gap}}),
\end{equation}
Since Softplus is positive and increasing, $w_{u,\text{pen}}>1$ and increases monotonically with $\Delta_{u,\mathrm{gap}}$.
Multiplying each user's negative term in Eq.~\ref{eq:sft_loss} by $w_{u,\text{pen}}$ gives
\begin{equation}
\label{eq:wpcl}
    \mathcal{L}_{\text{WPCL}} = -\mathop{\mathbb{E}}\limits_{u \in \mathcal{B}} \log \frac{e^{\text{s}(\mathbf{e}_u, \mathbf{e}_{i^+}) / \tau_{\text{WPCL}}}}{e^{\text{s}(\mathbf{e}_u, \mathbf{e}_{i^+}) / \tau_{\text{WPCL}}} + w_{u,\text{pen}} \sum_{i^- \in \mathcal{N}_u} e^{\text{s}(\mathbf{e}_u, \mathbf{e}_{i^-}) / \tau_{\text{WPCL}}}},
\end{equation}
\noindent where $\tau_{\text{WPCL}}$ is the temperature.
% Minimizing Eq.~\ref{eq:wpcl} also shrinks $\Delta_{u,\mathrm{gap}}$ by raising $\mathcal{M}_{u,\text{weak}}$, because $w_{u,\text{pen}}$ is not detached (Sec.~\ref{sec:wpcl_theory}).
Minimizing Eq.~\ref{eq:wpcl} also shrinks $\Delta_{u,\mathrm{gap}}$ by raising $\mathcal{M}_{u,\text{weak}}$, because $w_{u,\text{pen}}$ is not detached.


% Theoretical analysis removed from the paper (2026-10-26; kept for reference).
% \subsubsection{\textbf{Theoretical Analysis}}
% \label{sec:wpcl_theory}
% Through $w_{u,\text{pen}}$, $\mathcal{L}_{\text{WPCL}}$ raises the weak modality's margin and holds users with larger modality gaps to a higher standard, a wider score margin at the same loss (Proposition~\ref{prop:wpcl}; proof in Appendix~\ref{app:proof_wpcl}).

% \begin{proposition}[Weak-modality correction and margin offset]
% \label{prop:wpcl}
% Fix the strong/weak assignment, and differentiate the per-user term $\ell_u$ of Eq.~\ref{eq:wpcl} with respect to the VLM/LoRA parameters $\theta$.
% % (i) Split $\nabla_\theta\ell_u$ into the part with $w_{u,\text{pen}}$ held fixed and the part through $w_{u,\text{pen}}$, denoted $\nabla_\theta\ell_u^{\mathrm{gap}}$. The latter satisfies
% (i) The part of $\nabla_\theta\ell_u$ that flows through $w_{u,\text{pen}}$, denoted $\nabla_\theta\ell_u^{\mathrm{gap}}$, satisfies
% \begin{equation}
% \nabla_\theta \ell_u^{\mathrm{gap}} = -\gamma_u \nabla_\theta \mathcal{M}_{u,\text{weak}}, \qquad \gamma_u>0,
% \label{eq:wpcl_gradient_theory}
% \end{equation}
% so $-\nabla_\theta\ell_u^{\mathrm{gap}}$ is an ascent direction for $\mathcal{M}_{u,\text{weak}}$ whenever $\nabla_\theta\mathcal{M}_{u,\text{weak}}\neq\mathbf{0}$.
% (ii) Let $s_u^+=\tau_{\text{WPCL}}\log A_u$ and $r_u^-=\tau_{\text{WPCL}}\log B_u$, where $A_u$ and $B_u$ are the positive and summed negative exponentials in Eq.~\ref{eq:wpcl}. 
% To reach any fixed loss value, $\ell_u$ needs a score margin $s_u^+-r_u^-$ larger by $\tau_{\text{WPCL}}\log w_{u,\text{pen}}$ than $\mathcal{L}_{\text{SFT}}$ ($w_{u,\text{pen}}=1$) needs, and this offset increases with $\Delta_{u,\mathrm{gap}}$.
% \end{proposition}
% \noindent The gradient through $w_{u,\text{pen}}$ pushes the weak modality's positive score up and its negative scores down (Proposition~\ref{prop:wpcl}(i)): with $s_{u,j}^{w}=\text{s}(\mathbf e_u^{\mathrm{weak}},\mathbf e_j^{\mathrm{weak}})$ computed in $u$'s weak modality,
% \begin{equation}
% \frac{\partial\ell_u^{\mathrm{gap}}}{\partial s_{u,i^+}^{w}}=-\gamma_u,\quad
% \frac{\partial\ell_u^{\mathrm{gap}}}{\partial s_{u,i^-}^{w}}=\frac{\gamma_u}{|\mathcal N_u|},\; i^-\in\mathcal N_u.
% \label{eq:wpcl_score_gradient}
% \end{equation}

% \subsection{Training Objectives}

\subsection{Cross-modal Relational Topology Regularization}
\label{sec:crtr}

While $\mathcal{L}_{\text{WPCL}}$ strengthens the weak modality's discrimination, it does not relate the two modalities' similarity structures.
To satisfy the cross-modal consistency that multimodal SR also requires, we introduce \textit{Cross-modal Relational Topology Regularization} ($\mathcal{L}_{\text{CRTR}}$).
It softly aligns each user's relative similarity structure over the same items in the two modality spaces, so that fusing them improves recommendation.
Unlike direct alignment of $\mathbf{e}_u^T$ with $\mathbf{e}_u^V$, it does not require identical text and vision embeddings (Fig.~\ref{fig:method}).
$\mathcal{L}_{\text{WPCL}}$, which scores the summed embeddings, keeps the two spaces compatible for fusion, since $\mathcal{L}_{\text{CRTR}}$ ignores how they are oriented.

\subsubsection{\textbf{Modality-specific Relational Distribution}}
For user $u$, we define the modality-specific relational distribution $\mathbf{P}^m_u \in \mathbb{R}^{|\mathcal{C}|}$ over the candidate items $\mathcal{C}$, where each entry $\mathbf{P}_{u,j}^m$ is the softmax-normalized similarity between the sequence representation $\mathbf{e}_u^m$ of user $u$ and the candidate item embedding $\mathbf{e}_j^m$ of item $j$:
\begin{equation}
    \mathbf{P}_{u,j}^m = \frac{e^{\text{s}(\mathbf{e}_u^m, \mathbf{e}_j^m) / \tau_{\text{CRTR}}}}{\sum_{k\in\mathcal{C}} e^{\text{s}(\mathbf{e}_u^m, \mathbf{e}_k^m) / \tau_{\text{CRTR}}}},
\end{equation}
where $\tau_{\text{CRTR}}$ is a temperature parameter.
$\mathbf{P}_u^m$ captures the neighborhood of user $u$'s history among the candidate items in the modality-specific representation space.

\subsubsection{\textbf{Bidirectional Relational Consistency}}
For each user $u$, we then align the two neighborhoods $\mathbf{P}^T_u$ and $\mathbf{P}^V_u$ by minimizing their bidirectional Kullback--Leibler (KL) divergence:
\begin{equation}
\mathcal{L}_{\text{CRTR}} = \frac{1}{2}\mathop{\mathbb{E}}\limits_{u \in \mathcal{B}} \left( \mathrm{KL}(\mathbf{P}_u^T \| \mathbf{P}_u^V) + \mathrm{KL}(\mathbf{P}_u^V \| \mathbf{P}_u^T) \right).
\end{equation}
The bidirectional form pulls each modality's neighborhood toward the other's, rather than using the dominant modality as a fixed target.

% Theoretical analysis removed from the paper (2026-10-26; kept for reference).
% \subsubsection{\textbf{Theoretical Analysis}}
% \label{sec:crtr_theory}
% $\mathcal{L}_{\text{CRTR}}$ aligns each user's relative similarity structure in the two modality spaces \textbf{without forcing the text and vision embeddings to coincide} (Proposition~\ref{prop:crtr}; proof in Appendix~\ref{app:proof_crtr}).

% \begin{proposition}[Relational alignment without equal embeddings]
% \label{prop:crtr}
% Let both modalities compute $\mathbf P^m_u$ over a finite, nonempty item set $\mathcal{C}$ with the same $\tau_{\text{CRTR}}$ and cosine similarity $s^m_{u,j}=\text{s}(\mathbf e^m_u,\mathbf e^m_j)$. Let $D_u$ be the sum of the two KL terms of user $u$ in $\mathcal{L}_{\text{CRTR}}$.
% (i) $D_u=0$ if and only if $s^T_{u,j}-s^T_{u,k}=s^V_{u,j}-s^V_{u,k}$ for all $j,k\in\mathcal{C}$.
% (ii) $D_u$ is unchanged when one orthogonal map is applied to all embeddings of a modality (sequences and items), with independent maps for text and vision; hence $D_u=0$ does not require $\mathbf e^T_u=\mathbf e^V_u$.
% (iii) $\|\mathbf P_u^T-\mathbf P_u^V\|_1\le\sqrt{D_u}$.
% \end{proposition}

% \noindent By (iii), a small $D_u$ keeps the two neighborhoods close, but only among the in-batch candidates $\mathcal{C}$, not over the whole item set $\mathcal{I}$.

\subsection{Final Objective}
\label{sec:final_objective}
The final objective function $\mathcal{L}$ is defined as follows:
\begin{equation}
\mathcal{L} = \mathcal{L}_{\text{WPCL}} + \lambda \cdot \mathcal{L}_{\text{CRTR}},
\end{equation}
% where $\lambda$ balances discrimination and structural consistency.
where $\lambda$ balances discrimination and cross-modal consistency.
After training, VLM2Rec encodes a user's history into the fused sequence embedding $\mathbf{e}_u$ and directly ranks items by $\text{s}(\mathbf{e}_u,\mathbf{e}_i)$.

\begin{algorithm}[!h]
\caption{Training and inference of VLM2Rec and $\text{VLM2Rec}_{\text{plug}}$}
\label{alg:vlm2rec_pipeline}
\small
\begin{algorithmic}[1]
\Require Sequence--target pairs $(S_u,i^+)$, item catalog $\mathcal I$, VLM $\Phi$
\For{each minibatch $\mathcal B$ during training}
    \State Encode histories and targets with $P_T,P_V$; normalize last-token states and sum modalities.
    \State Compute per-modality margins; identify weak/strong, then $\Delta_{u,\mathrm{gap}}$ and $w_{u,\text{pen}}$ (Eqs.~\ref{eq:modality_gap} and~\ref{eq:wpen}).
    \State Compute $\mathcal L_{\text{WPCL}}$ (Eq.~\ref{eq:wpcl}) and $\mathcal L_{\text{CRTR}}$ from $\mathbf P_u^T$ and $\mathbf P_u^V$.
    \State Update VLM LoRA with $\mathcal L_{\text{WPCL}}+\lambda\mathcal L_{\text{CRTR}}$.
\EndFor
\State \textbf{Precompute:} Encode each item once into the fused item table $\mathbf E$.
\State \textbf{VLM2Rec:} At each request, encode the history with $\Phi$; rank items by $\text{s}(\mathbf e_u,\mathbf e_i)$ over $\mathbf E$.
\State \textbf{$\text{VLM2Rec}_{\text{plug}}$ (offline):} Replace the item embedding table of an SR backbone $f_\psi$ with the frozen $\mathbf E$ through $\mathbf W$; train only $\mathbf W$ and $\psi$.
\State \textbf{$\text{VLM2Rec}_{\text{plug}}$ (online):} Rank items with the trained SR backbone on $\mathbf E$, without VLM calls.
\end{algorithmic}
\end{algorithm}

% !TeX root = ../main.tex
\setlength{\fboxsep}{1pt}

\begin{table*}[!t]
\centering
\caption{Performance of VLM2Rec and $\text{VLM2Rec}_{\text{plug}}$. The best and second-best results within each setting are highlighted by \protect\colorbox{bestcolor}{\textbf{boldface}} and \protect\colorbox{secondcolor}{\underline{underline}}. * indicates statistical significance against the strongest baseline ($p < 0.05$, paired t-test over 5 seeds).}
\setlength{\tabcolsep}{1.5pt}
\resizebox{\linewidth}{!}{
\renewcommand{\arraystretch}{0.8}
\begin{tabular}{c|cccc|cccc|cccc|cccc|cccc}
\toprule
\multirow{2.5}{*}{\textbf{Method}} & \multicolumn{4}{c|}{\textbf{Toys}} & \multicolumn{4}{c|}{\textbf{Beauty}} & \multicolumn{4}{c|}{\textbf{Clothing}} & \multicolumn{4}{c|}{\textbf{Sports}} & \multicolumn{4}{c}{\textbf{MovieLens}} \\
\cmidrule(lr){2-5} \cmidrule(lr){6-9} \cmidrule(lr){10-13} \cmidrule(lr){14-17} \cmidrule(lr){18-21}
 & H@10 & H@20 & N@10 & N@20 & H@10 & H@20 & N@10 & N@20 & H@10 & H@20 & N@10 & N@20 & H@10 & H@20 & N@10 & N@20 & H@10 & H@20 & N@10 & N@20 \\
\midrule[\heavyrulewidth]
% ID-only results: arXiv:2603.17450v2, Table 10 (GRU4Rec/FMLP) and Table 3 (DiffuRec), Base columns.
GRU4Rec & 0.0662 & 0.0963 & 0.0362 & 0.0436 & 0.0638 & 0.0955 & 0.0350 & 0.0428 & 0.0312 & 0.0519 & 0.0153 & 0.0203 & 0.0326 & 0.0504 & 0.0171 & 0.0214 & 0.0962 & 0.1495 & 0.0471 & 0.0593 \\
FMLP & 0.0840 & 0.1138 & 0.0487 & 0.0561 & 0.0665 & 0.0987 & 0.0369 & 0.0451 & 0.0349 & 0.0564 & 0.0173 & 0.0224 & 0.0375 & 0.0564 & 0.0208 & 0.0255 & 0.1091 & 0.1602 & 0.0547 & 0.0673 \\
DiffuRec & 0.0890 & 0.1240 & 0.0500 & 0.0590 & 0.0851 & 0.1225 & 0.0471 & 0.0568 & 0.0488 & 0.0730 & 0.0252 & 0.0315 & 0.0473 & 0.0696 & 0.0250 & 0.0309 & 0.1296 & 0.1901 & 0.0643 & 0.0787 \\
SASRec & 0.0904 & 0.1221 & 0.0524 & 0.0602 & 0.0713 & 0.1053 & 0.0397 & 0.0483 & 0.0402 & 0.0648 & 0.0198 & 0.0260 & 0.0428 & 0.0648 & 0.0240 & 0.0294 & 0.1263 & 0.1830 & 0.0629 & 0.0773 \\
\midrule[\heavyrulewidth]
\multicolumn{21}{c}{\cellcolor{gray!30}\textbf{VLM2Rec Setting (without SR Backbone)}} \\ \midrule
BERT & 0.0550 & 0.0810 & 0.0310 & 0.0376 & 0.0433 & 0.0686 & 0.0223 & 0.0286 & 0.0234 & 0.0408 & 0.0113 & 0.0155 & 0.0244 & 0.0398 & 0.0119 & 0.0157 & 0.0679 & 0.1036 & 0.0320 & 0.0408 \\
SLIM & 0.0576 & 0.0850 & 0.0323 & 0.0390 & 0.0443 & 0.0708 & 0.0232 & 0.0297 & 0.0241 & 0.0426 & 0.0118 & 0.0162 & 0.0254 & 0.0412 & 0.0124 & 0.0164 & 0.0695 & 0.1066 & 0.0327 & 0.0420 \\
SLIM$^{+}$ & 0.0971 & 0.1379 & 0.0547 & 0.0647 & 0.0795 & 0.1213 & 0.0439 & 0.0544 & 0.0430 & 0.0673 & 0.0215 & 0.0276 & 0.0419 & 0.0640 & 0.0220 & 0.0276 & 0.1270 & 0.1909 & 0.0618 & 0.0780 \\
$\text{LLM}_{\text{Vanilla}}$ & 0.0781 & 0.1095 & 0.0436 & 0.0516 & 0.0593 & 0.0913 & 0.0342 & 0.0422 & 0.0332 & 0.0545 & 0.0169 & 0.0222 & 0.0351 & 0.0563 & 0.0192 & 0.0244 & 0.1133 & 0.1739 & 0.0547 & 0.0699 \\
$\text{LLM}_{\text{SFT}}$ & 0.1041 & 0.1467 & 0.0586 & 0.0693 & 0.0952 & 0.1539 & 0.0485 & 0.0631 & 0.0489 & 0.0745 & 0.0242 & 0.0306 & 0.0511 & 0.0756 & 0.0266 & 0.0327 & 0.1331 & 0.1975 & 0.0656 & 0.0817 \\
\midrule
$\text{VLM}_{\text{Vanilla}}$ & 0.0843 & 0.1197 & 0.0481 & 0.0570 & 0.0674 & 0.1054 & 0.0351 & 0.0442 & 0.0384 & 0.0610 & 0.0189 & 0.0245 & 0.0415 & 0.0605 & 0.0224 & 0.0271 & 0.1188 & 0.1740 & 0.0570 & 0.0711 \\
$\text{VLM}_{\text{SFT}}$ & \secondscore{0.1132} & 0.1456 & \secondscore{0.0652} & \secondscore{0.0732} & 0.0942 & 0.1527 & 0.0547 & 0.0695 & 0.0471 & \secondscore{0.0807} & 0.0241 & 0.0325 & 0.0506 & \secondscore{0.0817} & 0.0279 & \secondscore{0.0358} & 0.1671 & \secondscore{0.2527} & 0.0869 & \secondscore{0.1087} \\
MLLM-MSR & 0.1092 & 0.1504 & 0.0608 & 0.0711 & 0.0994 & 0.1573 & 0.0512 & 0.0655 & 0.0491 & 0.0750 & 0.0248 & 0.0313 & 0.0521 & 0.0788 & 0.0273 & 0.0342 & 0.1552 & 0.2259 & 0.0805 & 0.0982 \\
I-LLMRec & 0.1116 & \secondscore{0.1549} & 0.0621 & 0.0730 & \secondscore{0.1047} & \secondscore{0.1629} & \secondscore{0.0554} & \secondscore{0.0700} & \secondscore{0.0517} & 0.0803 & \secondscore{0.0258} & \secondscore{0.0331} & \secondscore{0.0545} & 0.0808 & \secondscore{0.0292} & 0.0357 & \secondscore{0.1699} & 0.2474 & \secondscore{0.0886} & 0.1085 \\
\midrule[\heavyrulewidth]
\textbf{VLM2Rec} & \bestscore{0.1254*} & \bestscore{0.1682*} & \bestscore{0.0720*} & \bestscore{0.0823*} & \bestscore{0.1161*} & \bestscore{0.1732*} & \bestscore{0.0641*} & \bestscore{0.0784*} & \bestscore{0.0553*} & \bestscore{0.0875*} & \bestscore{0.0290*} & \bestscore{0.0370*} & \bestscore{0.0585*} & \bestscore{0.0904*} & \bestscore{0.0316*} & \bestscore{0.0396*} & \bestscore{0.1899*} & \bestscore{0.2783*} & \bestscore{0.1029*} & \bestscore{0.1252*} \\
\textit{Imp. (\%)} & 10.78 & 8.59 & 10.43 & 12.43 & 10.89 & 6.32 & 15.70 & 12.00 & 6.96 & 8.43 & 12.40 & 11.78 & 7.34 & 10.65 & 8.22 & 10.61 & 11.77 & 10.13 & 16.14 & 15.18 \\
\midrule[\heavyrulewidth]
\multicolumn{21}{c}{\cellcolor{gray!30}\textbf{$\text{VLM2Rec}_{\text{plug}}$ Setting (Backbone: SASRec)}} \\ \midrule
BERT & 0.0702 & 0.1037 & 0.0387 & 0.0471 & 0.0601 & 0.0952 & 0.0308 & 0.0395 & 0.0303 & 0.0517 & 0.0144 & 0.0197 & 0.0314 & 0.0511 & 0.0157 & 0.0207 & 0.0884 & 0.1337 & 0.0397 & 0.0510 \\
SLIM & 0.0733 & 0.1077 & 0.0412 & 0.0497 & 0.0606 & 0.0957 & 0.0311 & 0.0397 & 0.0306 & 0.0525 & 0.0147 & 0.0201 & 0.0324 & 0.0519 & 0.0162 & 0.0210 & 0.0901 & 0.1347 & 0.0403 & 0.0512 \\
SLIM$^{+}$ & 0.0961 & 0.1332 & 0.0561 & 0.0655 & 0.0812 & 0.1205 & 0.0443 & 0.0540 & 0.0434 & 0.0681 & 0.0218 & 0.0280 & 0.0427 & 0.0656 & 0.0229 & 0.0287 & 0.1555 & 0.2341 & 0.0779 & 0.0978 \\
LLMEmb & 0.0985 & 0.1371 & 0.0574 & 0.0667 & 0.0854 & 0.1247 & 0.0468 & 0.0564 & 0.0452 & 0.0692 & 0.0234 & 0.0292 & 0.0442 & 0.0680 & 0.0241 & 0.0299 & 0.1588 & 0.2356 & 0.0796 & 0.0985 \\
LLM2Rec & 0.1004 & 0.1403 & 0.0592 & 0.0691 & 0.0896 & 0.1299 & 0.0493 & 0.0594 & 0.0470 & 0.0718 & 0.0244 & 0.0305 & 0.0461 & 0.0691 & 0.0253 & 0.0312 & 0.1598 & 0.2367 & 0.0827 & 0.1018 \\
$\text{LLM}_{\text{Vanilla}}$ & 0.0816 & 0.1142 & 0.0457 & 0.0540 & 0.0645 & 0.0985 & 0.0337 & 0.0425 & 0.0341 & 0.0563 & 0.0164 & 0.0220 & 0.0367 & 0.0575 & 0.0199 & 0.0250 & 0.1095 & 0.1654 & 0.0538 & 0.0678 \\
$\text{LLM}_{\text{SFT}}$ & 0.1015 & 0.1407 & 0.0616 & 0.0715 & 0.0948 & 0.1365 & 0.0506 & 0.0612 & 0.0479 & 0.0724 & 0.0249 & 0.0310 & 0.0492 & 0.0716 & 0.0268 & 0.0325 & 0.1624 & 0.2374 & 0.0865 & 0.1052 \\
\midrule
CLIP & 0.0735 & 0.1044 & 0.0408 & 0.0484 & 0.0663 & 0.0997 & 0.0358 & 0.0440 & 0.0344 & 0.0578 & 0.0173 & 0.0233 & 0.0357 & 0.0563 & 0.0191 & 0.0243 & 0.1033 & 0.1560 & 0.0506 & 0.0636 \\
$\text{VLM}_{\text{Prompt}}$ & 0.0792 & 0.1123 & 0.0449 & 0.0530 & 0.0679 & 0.1008 & 0.0367 & 0.0449 & 0.0359 & 0.0594 & 0.0178 & 0.0235 & 0.0374 & 0.0588 & 0.0204 & 0.0256 & 0.1115 & 0.1668 & 0.0548 & 0.0690 \\
NoteLLM-2 & 0.1049 & 0.1438 & 0.0596 & 0.0690 & 0.0903 & 0.1309 & 0.0507 & 0.0608 & 0.0492 & 0.0744 & 0.0258 & 0.0320 & 0.0470 & 0.0730 & 0.0261 & 0.0324 & 0.1631 & 0.2396 & 0.0885 & 0.1072 \\
$\text{VLM}_{\text{Vanilla}}$ & 0.0841 & 0.1178 & 0.0491 & 0.0577 & 0.0685 & 0.1031 & 0.0374 & 0.0463 & 0.0365 & 0.0596 & 0.0183 & 0.0240 & 0.0383 & 0.0601 & 0.0208 & 0.0262 & 0.1152 & 0.1695 & 0.0559 & 0.0695 \\
$\text{VLM}_{\text{SFT}}$ & 0.1058 & \secondscore{0.1461} & 0.0598 & 0.0699 & 0.0922 & 0.1349 & 0.0515 & 0.0622 & \secondscore{0.0499} & 0.0755 & \secondscore{0.0260} & 0.0325 & 0.0469 & \secondscore{0.0738} & 0.0260 & 0.0326 & \secondscore{0.1649} & 0.2365 & \secondscore{0.0902} & 0.1082 \\
MLLM-MSR & 0.1038 & 0.1421 & 0.0611 & 0.0706 & 0.0940 & 0.1340 & 0.0506 & 0.0604 & 0.0482 & 0.0730 & 0.0253 & 0.0316 & 0.0494 & 0.0725 & 0.0272 & 0.0329 & 0.1627 & 0.2402 & 0.0864 & 0.1057 \\
I-LLMRec & \secondscore{0.1064} & 0.1442 & \secondscore{0.0624} & \secondscore{0.0717} & \secondscore{0.0967} & \secondscore{0.1390} & \secondscore{0.0539} & \secondscore{0.0644} & 0.0483 & \secondscore{0.0759} & 0.0258 & \secondscore{0.0328} & \secondscore{0.0503} & 0.0737 & \secondscore{0.0278} & \secondscore{0.0337} & 0.1644 & \secondscore{0.2461} & 0.0881 & \secondscore{0.1085} \\
\midrule[\heavyrulewidth]
\textbf{$\text{VLM2Rec}_{\text{plug}}$} & \bestscore{0.1169*} & \bestscore{0.1588*} & \bestscore{0.0684*} & \bestscore{0.0786*} & \bestscore{0.1080*} & \bestscore{0.1511*} & \bestscore{0.0618*} & \bestscore{0.0728*} & \bestscore{0.0540*} & \bestscore{0.0823*} & \bestscore{0.0276*} & \bestscore{0.0345*} & \bestscore{0.0566*} & \bestscore{0.0849*} & \bestscore{0.0307*} & \bestscore{0.0378*} & \bestscore{0.1782*} & \bestscore{0.2718*} & \bestscore{0.0977*} & \bestscore{0.1216*} \\
% \midrule
\textit{Imp. (\%)} & 9.87 & 8.69 & 9.62 & 9.62 & 11.69 & 8.71 & 14.66 & 13.04 & 8.22 & 8.43 & 6.15 & 5.18 & 12.52 & 15.04 & 10.43 & 12.17 & 8.07 & 10.44 & 8.31 & 12.07 \\
\bottomrule
\end{tabular}
}
\label{tab:perf_task2}
\end{table*}



\subsection{$\text{VLM2Rec}_{\text{plug}}$: Plugging into SR Backbones}
\label{sec:deployment}
VLM2Rec runs the VLM on each user's latest history at every request, which is far heavier than a typical SR model.
% Item embeddings, however, depend only on the items and can be computed once offline.
However, item embeddings depend only on the items and can be computed once offline.
% In this setting, we encode every item once into a frozen item table $\mathbf{E}=\{\mathbf{e}_i\}_{i\in\mathcal{I}}$ with the trained VLM2Rec model.
We therefore propose the $\text{VLM2Rec}_{\text{plug}}$ setting, in which we encode every item with the trained VLM2Rec model into a frozen item table $\mathbf{E}=\{\mathbf{e}_i\}_{i\in\mathcal{I}}$ (Fig.~\ref{fig:method}).
$\mathbf{E}$ inherits the sequence-aware knowledge that the model learned by reasoning over multimodal histories, even though each item is encoded alone.
This table then replaces the item embedding table of an SR backbone $f_\psi$ through a linear adapter $\mathbf{W}$:
\begin{equation}
\hat{y}_{u,i}=f_\psi(\mathbf{W}\mathbf{e}_{i_1},\dots,\mathbf{W}\mathbf{e}_{i_L})^{\top}\mathbf{W}\mathbf{e}_i.
\label{eq:sr_score}
\end{equation}
Only $\mathbf{W}$ and $\psi$ are trained with the backbone's loss, so the backbone learns item transitions on embeddings that already carry multimodal semantics and sequence-level CF.
The serving latency matches that of the SR backbone alone, because $\mathbf{E}$ is precomputed and no VLM runs per request.
Algorithm~\ref{alg:vlm2rec_pipeline} summarizes training and inference of both settings.
